Few-shot prompts are often assembled by hand, so their demonstrations can carry spurious statistics: one label may dominate, sit at the end of the prompt, or come with longer inputs. Studies of GPT-3-era base models found that such statistics pull predictions strongly toward the favoured label. We ask how large these effects are in current instruction-tuned models and how they change with scale. We build a fully paired benchmark of 10,000 prompts per model over four classification tasks. It covers four manipulations of a 10-shot demonstration set: frequency skew (80/20), recency, primacy, and label--length correlation. The favoured label is rotated through all labels, so intrinsic label preferences cancel. We score eight local models: Qwen2.5-Instruct from 0.5B to 14B, a Qwen2.5-1.5B base model, and two Llama-3 models. Every bias is far smaller than hypothesised. Frequency skew is the largest effect in small models (+4.3 to +16.9 percentage points) and falls to about zero at 8--14B ($-3.7$ points per tenfold increase in parameters). Recency does not pull predictions toward the last label; from 3B upward it produces a small, significant \emph{anti}-recency effect ($-1.7$ to $-3.6$ points). Label--length correlation has no reliable effect, and calibration error does not change reliably. Large shifts appear only under extreme manipulations, such as all demonstrations sharing one label (up to +36 points), and mostly on the ambiguous paraphrase task. For models of 3B and above, re-drawing a balanced demonstration set flips more predictions (4--8\%) than any systematic bias we measure.
